% subsection: 対数型 \, $a_{n+1}=pa_n^q$

\subsection{対数型 \, $a_{n+1}=pa_n^q$}
  今までは$a_n$に対する操作は定数倍したり関数が掛けられたりするだけだった.
  しかし,\,$a_n$に指数がついたらどうだろうか.
  この場合はかなり難しい発想の変形が必要である.
  それは対数を取るという操作である.
  実際にやってみよう.
  \[
    a_{n+1}=pa_n^q
  \]
  ただし,対数を取るために$a_n>0$とし,底に選ぶ$p$も$p>0,\,p\ne1$を満たすとする.
  この式について,両辺の対数を取ると
  \begin{align*}
    \log_p{a_{n+1}}&=\log_p{pa_n^q}\\
    &=\log_p{p}+\log_p{a_n^q}\\
    &=q\log_p{a_n}+1
  \end{align*}
  $b_n=\log_p{a_n}$とおくと,
  \begin{align*}
    b_{n+1}=qb_n+1
  \end{align*}
  あとはただの特性方程式型の問題である.
  底をいくらにするかは自由だが$p$に関連する数だと扱いやすい.
  たとえば$p=k^m$の時は$k$でとっておくことが多い.
  作問者側は解答がシンプルな形になるようにするため,より小さい底を取ることで答えが求めやすくなる.

  また,対数の強みは積を和に,商を差にできるところにある.
  漸化式が全て積や商のみで表される場合にその強みが現れる.
  これはここで解説するよりも演習問題を通して経験するほうがわかりやすいと判断したため,この内容は(2)で取り扱っている.

  \noindent\textbf{【例題】}
  \begin{enumerate}[label=(\arabic*),parsep=3pt]
    \item $\displaystyle a_1=2,\quad a_{n+1}=2a_n^2$
    \item $\displaystyle a_1=\frac{1}{2},\quad a_{n+1}=\frac{1}{a_n}$
    \item $\displaystyle a_1=1,\quad a_{n+1}=2^na_n$
    \item $\displaystyle a_1=5,\quad a_{n+1}=a_n^2-2a_n+2$
    \item $\displaystyle a_1=2,\quad a_{n+1}=\frac{a_n^2}{2^n}$
    \item $\displaystyle a_1=1,\quad a_{n+1}=(n+1)a_n$
  \end{enumerate}

  \noindent\textbf{【方針】}
  \begin{enumerate}[label=(\arabic*),parsep=3pt]
    \item 両辺の底を$2$とする対数を取り,\,$b_n=\log_2a_n$とおいて一次の漸化式に変形します.その結果,漸化式は特性方程式型になります.
    \item 一見すると分数型(後述)のように見えますが,対数型として解くことができます.
    \item $f(n)=2^n$が指数関数なので,対数を取ると和の形に帰着できます.
    \item うまく変形すると対数型になります.
    \item 両辺の対数を取ると,二次の漸化式が一次の漸化式に変わり,さらに等差数列の和を使って解くことができます.
    \item 階比型に見えますが,対数型としても解けます.
  \end{enumerate}

  \noindent\textbf{【解答】}
  \begin{enumerate}[label=(\arabic*),parsep=3pt]
    \item $\displaystyle a_n=2^{2^{n}-1}$
    \item $a_n=2^{(-1)^{n}}$
    \item $a_n=2^{\frac{n(n-1)}{2}}$
    \item $a_n=2^{2^n}+1$
    \item $\displaystyle a_n=2^{n-1}-n+1$
    \item $a_n=n!$
  \end{enumerate}

  \noindent\textbf{【解法】}
  \begin{enumerate}[label=(\arabic*),parsep=3pt]
    \item $\displaystyle a_1=2,\quad a_{n+1}=2a_n^2$\\
          両辺の底を$2$とする対数を取ると,
          \begin{align*}
            \log_2a_{n+1}
              &=\log_2{2a_n^2}\\
              &=\log_2{2}+\log_2{a_n^2}\\
              &=2\log_2{a_n}+1
          \end{align*}
          ここで
          \[
            b_n=\log_2a_n
          \]
          とおくと,
          \begin{align*}
            &b_{n+1}=2b_n+1\\
            &b_{n+1}+1=2(b_n+1)
          \end{align*}
          ここで
          \[
            c_n=b_n+1
          \]
          とおくと,
          \[
            c_{n+1}=2c_n
          \]
          $c_1$は,
          \[
            c_1=b_1+1=\log_2{a_1}+1=\log_2{2}+1=1+1=2
          \]
          従って,
          \begin{align*}
            &c_n=2\cdot2^{n-1}=2^n\\
            &b_n+1=2^n\\
            &b_n=2^n-1\\
            &a_n=2^{2^{n}-1}
          \end{align*}
    \item $\displaystyle a_1=\frac12,\quad a_{n+1}=\frac{1}{a_n}$\\
          与式の両辺の底を2とする対数を取ると,
          \begin{align*}
            &\log_2{a_{n+1}}=\log_2{\frac{1}{a_n}}\\
            &\log_2{a_{n+1}}=\log_2{1}-\log_2{a_n}\\
            &\log_2{a_{n+1}}=-\log_2{a_n}\\
            &\log_2{a_n}=\log_2{a_1}(-1)^{n-1}\\
            &\log_2{a_n}=\log_2{\frac{1}{2}}(-1)^{n-1}\\
            &\log_2{a_n}=(-1)^{n}\\
            &\therefore \quad a_n=2^{(-1)^{n}}
          \end{align*}
    \item $\displaystyle a_1=1,\quad a_{n+1}=2^na_n$\\
          両辺の自然対数を取ると
          \[
            \log a_{n+1}=n\log2+\log a_n
          \]
          したがって
          \begin{align*}
            \log a_n&=\log a_1+\sum_{k=1}^{n-1} k\log2\\
            &=\frac{(n-1)n}{2}\log2\\
            &\therefore \quad a_n=2^{\frac{n(n-1)}{2}}
          \end{align*}
    \item $\displaystyle a_1=5,\quad a_{n+1}=a_n^2-2a_n+2$\\
          与式を変形して
          \[
            a_{n+1}-1=(a_n-1)^2
          \]
          両辺の底を2とする対数を取ると,
          \begin{align*}
            &\log_2{(a_{n+1}-1)}=\log_2{(a_n-1)^2}\\
            &\log_2{(a_{n+1}-1)}=2\log_2{(a_n-1)}\\
            &\log_2{(a_n-1)}=2^{n-1}\cdot\log_2{(a_1-1)}\\
            &\log_2{(a_n-1)}=2^{n-1}\cdot\log_2{4}\\
            &\log_2{(a_n-1)}=2^{n}\\
            &a_n-1=2^{2^{n}}\\
            &\therefore \quad a_n=2^{2^{n}}+1
          \end{align*}
    \item $\displaystyle a_1=2,\quad a_{n+1}=\frac{a_n^2}{2^n}$\\
          両辺を底を2とする自然体数を取って
          \[
            equation
          \]
    \item $\displaystyle a_1=1,\quad a_{n+1}=(n+1)a_n$\\
          両辺の自然対数を取って,
          \[
            \log a_{n+1}=\log{n+1}+\log a_n
          \]
          したがって
          \begin{align*}
            \log a_n&=\log a_1 + \log{2}+\cdots+ \log{n}\\
            &=\log 1 + \log{2}+\cdots +\log{n}\\
            &=\log n!\\
            &\therefore \quad a_n=n!
          \end{align*}
  \end{enumerate}
  \noindent\textbf{【補足】}

  (3)では底を自然対数としたが,底は任意の値で構わない.

